\section{Further structure of the one-bit model}\label{sec:further-structure}

The one-bit model has two further structural properties: affine dependence on the initial table,
and a delay before an update rule becomes observable in the output.

\subsection{Affine dependence on the initial table}

The split in Section~\ref{sec:structure} concerns linearity in the \emph{successor sequence}.
With the source held fixed, a different fact holds: the transform is affine in its initial table
for all sixteen wirings, including those with $d_3=1$.

To see this, factor the ANF as
\begin{equation}\label{eq:factored-appendix}
  w(s,x)=\underbrace{(d_0+d_2x)}_{\alpha(x)}
        +\underbrace{(d_1+d_3x)}_{\beta(x)}s .
\end{equation}
Fix one context with successor sequence $y_0,y_1,\dots$ and initial cell $q$.  Its recurrence has
the form
\[
  u_{j+1}=\alpha(y_j)+\beta(y_j)u_j.
\]
Induction gives
\begin{equation}\label{eq:affine-in-q}
  u_j=\mu_j+\lambda_jq,
  \qquad
  \lambda_j=\prod_{i<j}\beta(y_i),
\end{equation}
where $\mu_j$ and $\lambda_j$ depend on the fixed source but not on $q$.  Since
$r_{k_j}=y_j+u_j$, every residual in that context is affine in its own initial cell.

Cells never interact, and the address sequence depends on the source and prefix rather than the
initial table.  Consequently, for fixed $(\pi,x)$ there is a linear map
$P_{\pi,x}:\F^{2^m}\to\F^n$ such that
\begin{equation}\label{eq:table-affine-appendix}
  \cxor_w(q,\pi,x)
  =\cxor_w(0,\pi,x)+P_{\pi,x}q
\end{equation}
for every wiring $w$.  This does not contradict the nonlinear classification: when $d_3=1$,
$\beta(x)$ is a source-dependent coefficient, which is constant only because $x$ is held fixed
in Equation~\eqref{eq:table-affine-appendix}.

The columns of $P_{\pi,x}$ have disjoint supports: changing cell $a$ can affect only visits to
context $a$.  Moreover, a visited cell always affects its first visit, because the empty product
in Equation~\eqref{eq:affine-in-q} is $\lambda_0=1$.  It follows immediately that
\begin{equation}\label{eq:rank-appendix}
  \operatorname{rank}P_{\pi,x}
  =\bigl|\{a:\text{$a$ is visited by $(\pi,x)$}\}\bigr|
  \le 2^m .
\end{equation}
Thus a known source/residual pair determines the initial bits of all visited cells, while
unvisited cells are necessarily unobservable.  Equation~\eqref{eq:delta-horizon} refines the same
calculation by describing the full support of each nonzero column.

\subsection{When does the wiring become visible?}

At a context's first visit, the emitted residual is $y_0+q$ and is independent of the update rule;
the rule acts only after that bit has been emitted.  Differences between two wirings can therefore
be observed only on a later visit to a cell on which their earlier updates disagreed.

$\W5$ and $\W6$ have an additional one-visit agreement when the table starts at zero.  On the
first visit both store $y_0$.  They consequently emit the same residual on the second visit, but
then store respectively
\[
  \W5(y_0,y_1)=y_1,
  \qquad
  \W6(y_0,y_1)=y_0+y_1.
\]
If $y_0=1$, that difference becomes observable on the third visit to the context.  For a zero
prefix, the tape $1\,0^{m+2}$ supplies such a trace in the all-zero context and separates the two
rules at its last symbol; its length is $m+3$.  The verification harness confirms that no shorter
tape separates them for $m\le6$.  Thus the $\W5$/$\W6$ distinction is invisible until some
context has accumulated enough visit history.
